\chapter*{Введение}
\label{sec:afterwords}
\addcontentsline{toc}{chapter}{Введение}

Повышение эффективности сельского хозяйства является одной из ключевых задач современного мира, особенно в условиях изменения климата и ограниченности ресурсов. В этом контексте оптимизация процессов выращивания сельскохозяйственных культур, в частности дынь, приобретает особую актуальность. Дыни, будучи ценной продовольственной культурой, подвержены множеству заболеваний и стрессовых факторов, которые могут существенно снизить урожайность и качество продукции. Своевременная и точная диагностика состояния растений, а также оперативное принятие мер по предотвращению и лечению заболеваний, являются критически важными для обеспечения стабильного урожая.

Традиционные методы мониторинга и диагностики, основанные на ручном осмотре полей и опыте агрономов, зачастую оказываются неэффективными при больших объемах производства. Они требуют значительных временных и трудовых затрат, а также высокой квалификации специалистов, что делает их дорогостоящими и подверженными человеческому фактору. Кроме того, многие заболевания проявляются на ранних стадиях неявными симптомами, которые трудно обнаружить без специальных средств, что приводит к задержкам в принятии решений и, как следствие, к потерям урожая.

В связи с этим возрастает потребность в разработке и внедрении интеллектуальных систем, способных автоматизировать процессы мониторинга, диагностики и принятия решений в сельском хозяйстве. Интегрированные экспертные системы (ИЭС) представляют собой мощный инструмент для решения подобных задач, позволяя аккумулировать и эффективно использовать знания экспертов, а также обрабатывать большие объемы данных с различных источников.

\textbf{Актуальность темы} обусловлена необходимостью повышения эффективности процессов диагностики заболеваний дынь в условиях нестабильного содержания питательных веществ в почве. Разработка динамической интегрированной экспертной системы позволит автоматизировать процесс принятия решений, сократить потери урожая и оптимизировать использование ресурсов в сельском хозяйстве. Нестабильность содержания питательных веществ в почве является одной из важнейших проблем современного земледелия, так как она напрямую влияет на развитие растений и их устойчивость к заболеваниям. При этом изменение концентрации основных макро- и микроэлементов может быть вызвано различными факторами, такими как погодные условия, интенсивность осадков, особенности состава почвы и применение агрохимикатов.

\textbf{Целью данной работы} является разработка прототипа динамической интегрированной экспертной системы для диагностики заболеваний дынь по визуальным признакам, выращиваемых в условиях нестабильного содержания питательных веществ в почве.

Для достижения поставленной цели необходимо решить следующие \textbf{задачи}:
\begin{enumerate}
    \item Проанализировать проблемную область, выявить основные источники знаний и методы их получения. Систематизировать информацию о заболеваниях дынь и дефицитах питательных веществ, их визуальных симптомах и влиянии условий внешней среды на их развитие.
    \item Построить фрагмент поля знаний для данной проблемной области, включающий правила вывода на естественном языке и правила на языке представления знаний системы G2.
    \item Разработать имитационную модель, описывающую динамику изменения состояния дынь под воздействием внешних факторов (климатические условия, состав почвы, питательные вещества, загрязняющие вещества).
    \item Разработать общую архитектуру и структуру прототипа интегрированной экспертной системы, определить ее основные компоненты и их взаимодействие.
    \item Реализовать прототип интегрированной экспертной системы, включая базу знаний и механизм вывода, используя инструментальное средство G2.
    \item Провести тестирование прототипа ИЭС в двух режимах функционирования: режиме консультации (проверка корректности диагностики и рекомендаций) и режиме имитации (проверка адекватности моделирования динамики состояния растений).
    \item Проанализировать результаты тестирования и оценить эффективность разработанной системы для практического применения в сельском хозяйстве.
\end{enumerate}

\textbf{Объектом исследования} является процесс диагностики заболеваний дынь по визуальным признакам в условиях нестабильного содержания питательных веществ в почве.

\textbf{Предметом исследования} являются методы и модели построения динамических интегрированных экспертных систем для поддержки принятия решений в сельском хозяйстве.

\textbf{Научная новизна} работы заключается в следующем:
\begin{itemize}
    \item Разработана теоретико-множественная модель имитационной модели, адаптированная для оценки динамики состояния дынь в условиях нестабильного содержания питательных веществ в почве, учитывающая влияние климатических факторов и загрязнений.
    \item Предложен набор правил для базы знаний динамической интегрированной экспертной системы, позволяющий диагностировать заболевания дынь и дефициты питательных веществ на основе визуальных признаков и параметров внешней среды.
    \item Разработаны тестовые сценарии для оценки адекватности и эффективности прототипа динамической интегрированной экспертной системы в режимах консультации и имитации, учитывающие динамику изменения параметров окружающей среды.
\end{itemize}

\textbf{Практическая значимость} работы состоит в создании прототипа динамической интегрированной экспертной системы, который может быть использован агрономами для оперативной диагностики заболеваний дынь, оптимизации агротехнических мероприятий и повышения урожайности. Система может быть интегрирована в существующие системы мониторинга сельского хозяйства, что позволит снизить потери от заболеваний и дефицитов питания, а также оптимизировать использование ресурсов.

\textbf{Структура работы} состоит из введения, четырех глав, заключения, списка литературы и приложений.
